\begin{frame}{Método de Euler}

\begin{block}{Problema}
Por el criterio de Weyl, para demostrar que una sucesión es
equidistribuida debemos estudiar sumas de la forma
\[\sum_{n=1}^{N} e^{2\pi i f(n)}.\]
\end{block}

\begin{alertblock}{Idea}
Transformar sumas discretas en integrales que puedan
estimarse con precisión.
\end{alertblock}
\vspace{0.3cm}
Herramienta principal:
\[
\boxed{\text{Fórmula de sumación de Euler}}
\]
\end{frame}

\begin{frame}{Fórmula de sumación de Euler}

\begin{block}{Proposición}
Sea $f:[1,N]\to\mathbb C$ con $f'\in C([1,N])$.

Entonces
\[
\sum_{n=1}^{N}f(n)=
\int_1^N f(x)\,dx + \frac{f(1)+f(N)}{2} +
\int_1^N\left(\langle x\rangle-\frac12\right)f'(x)\,dx.
\]
\end{block}

\begin{alertblock}{Interpretación}
La suma discreta se expresa mediante integrales
más fáciles de estimar.
\end{alertblock}

\end{frame}

\begin{frame}{Ejemplo: la sucesión $a\log n$}

\begin{block}{Resultado}
La sucesión $(a\log n)$ no es equidistribuida módulo $1$.
\end{block}

Aplicando la fórmula de Euler a
\[
f(x)=e^{2\pi i a\log x}
\]
obtenemos una estimación de la suma exponencial asociada.

\begin{alertblock}{Conclusión}
El criterio de Weyl falla para $k=1$.
Por tanto, $(a\log n)$ no es equidistribuida.
\end{alertblock}

\end{frame}

\begin{frame}{Ejemplo: la sucesión $an^\sigma$}

\begin{block}{Hipótesis}
\[
a\neq 0,
\qquad
0<\sigma<1.
\]
\end{block}

Aplicando nuevamente la fórmula de Euler a

\[
f(x)=e^{2\pi i a x^\sigma}
\]

se obtiene

\[
\sum_{n=1}^{N}
e^{2\pi i a n^\sigma}
=
o(N).
\]


\begin{alertblock}{Por el criterio de Weyl}
\[
(an^\sigma)
\]

es equidistribuida módulo $1$.
\end{alertblock}

\end{frame}

\begin{frame}{Hacia el teorema de Fejér}

Los ejemplos anteriores sugieren que la equidistribución
depende del comportamiento de la derivada.

\vspace{0.3cm}

\begin{center}

\[
g(x)=\log x
\]

\[
g(x)=x^\sigma
\]

\end{center}

\vspace{0.3cm}

\pause

\begin{alertblock}{Pregunta}
¿Qué condiciones sobre una función $g$
garantizan que $(g(n))$ sea equidistribuida?
\end{alertblock}

\end{frame}


\begin{frame}{Teorema de Fejér}

\begin{block}{Teorema}
Sea $g:[1,\infty)\to\mathbb R$ derivable.

Si

\[
g'(x)\to 0,
\]

\[
x|g'(x)|\to\infty,
\]

y $g'(x)$ es monótona para $x$ suficientemente grande,

entonces

\[
(g(n))
\]

es equidistribuida módulo $1$.
\end{block}

\end{frame}

\begin{frame}{Aplicaciones del teorema de Fejér}

\begin{block}{Consecuencias}
Las siguientes sucesiones son equidistribuidas módulo $1$:
\end{block}

\[
(an^\sigma),
\qquad
0<\sigma<1.
\]

\[
(an^\sigma \log^\tau n).
\]

\[
(a\log^\tau n),
\qquad
\tau>1.
\]

\[
(an\log^\tau n),
\qquad
\tau<0.
\]

\pause

\begin{alertblock}{Idea principal}
El método de Euler proporciona una herramienta
eficaz para estimar sumas exponenciales.
\end{alertblock}

\end{frame}